%---------------------------------------------------------------
% Kurzfassung
%---------------------------------------------------------------

\header{Kurzfassung}
Die Zementindustrie ist gegenwärtig für rund 7 \% der globalen CO\textsubscript{2}-Emissionen verantwortlich, wobei ein erheblicher Teil des elektrischen Energiebedarfs auf die Vermahlung entfällt. Vor diesem Hintergrund bietet die Betriebspunktoptimierung von Mahlanlagen ein signifikantes Potenzial zur Reduktion der spezifischen Energie, also der benötigten Energie pro Tonne produzierten Materials. Die Grundlage solcher Optimierungen bilden Simulationsmodelle, deren Performance jedoch im zeitlichen Verlauf aufgrund von Einflüssen wie Materialvariationen, Verschleiß und Wartungen beeinträchtigt wird. Im Rahmen der Untersuchung kann dieses Phänomen, aufgrund der Art der zugrunde liegenden Verschiebung der Datengrundlage als Real Concept Drift klassifiziert werden. \newline
Durch die künstliche Simulation von Concept Drift wird das identifizierte Problem anhand des Wiener Modells generalisiert und zusätzlich auf den Benchmark-Datensatz des Tennessee Eastman Process übertragen, der für diesen Zweck modifiziert wird. \newline
Kern der Arbeit ist die Untersuchung von Adaption der Modelle an veränderliche Bedingungen im industriellen Umfeld unter Verwendung neuronaler Netze, um die langfristige Simulationsgüte sicherzustellen. Dabei wird das Stability-Plasticity-Dilemma adressiert, also die Herausforderung des Modells, gleichzeitig neues Wissen aufzunehmen und vorhandenes Wissen zu erhalten. \newline
Die Ergebnisse zeigen eine signifikante Verbesserung der Modellperformance durch die Adaption an veränderte Prozessbedingungen. Die vorliegende Arbeit liefert damit einen Beitrag zur Schaffung robuster maschineller Lernmodelle für die industrielle Praxis und zeigt messbare Vorteile für die Energieeffizienz vergleichbarer Prozesse auf.

\cleardoublepage